\documentclass[10pt]{article}
\usepackage[utf8]{inputenc}
\usepackage[T5,T1]{fontenc}
\usepackage[vietnamese]{babel}
\usepackage{amsmath,amssymb,amsfonts}
\usepackage{graphicx}
\usepackage{xcolor}
\usepackage{tikz}
\usepackage{enumitem}
\usepackage{geometry}
\usepackage{fancyhdr}
\usepackage[skins,breakable]{tcolorbox}

% Thiết lập kích thước trang chuẩn Letter của giáo trình Stewart Calculus
\geometry{
    letterpaper,
    top=1.3cm,
    bottom=1.1cm,
    left=1.6cm,
    right=1.6cm,
    headheight=14pt,
    headsep=10pt
}

% Màu nhận diện Stewart Calculus 9th Edition
\definecolor{stewartcyan}{RGB}{0, 130, 155}       % Xanh cyan Mục 3.5, icon, số đề mục
\definecolor{projectgreen}{RGB}{0, 115, 75}       % Xanh lục đậm tiêu đề Dự án
\definecolor{projectbg}{RGB}{236, 245, 239}        % Nền xanh nhạt dịu của hộp Dự án
\definecolor{copyrightgray}{RGB}{125, 125, 125}   % Xám bản quyền chân trang

\pagestyle{fancy}
\fancyhf{}
\renewcommand{\headrulewidth}{0pt}
\fancyhead[R]{\fontsize{8.5pt}{10pt}\selectfont\textsf{\textbf{\color{stewartcyan}MỤC 3.5}}\hspace{1.5em}\textsf{Đạo hàm ẩn}\hspace{2.5em}\textsf{\textbf{209}}}

\setlength{\parindent}{0pt}

\begin{document}

% ==============================================================================
% PHẦN TRÊN: BÀI TẬP 99 - 102 (BỐ CỤC 2 CỘT)
% ==============================================================================
\noindent
\begin{minipage}[t]{0.46\textwidth}
    \vspace{0pt}
    \noindent\textbf{99.} Gọi $c$ là hoành độ giao điểm của tiếp tuyến với đường cong $y = b^x$ ($b > 0, b \ne 1$) tại điểm $(a, b^a)$ với trục hoành. Chứng minh rằng khoảng cách giữa hai điểm $(a, 0)$ và $(c, 0)$ là như nhau đối với mọi giá trị của $a$.

    \vspace{0.15cm}
    \begin{center}
        \includegraphics[width=0.62\linewidth]{images/p0244_fig1.png}
    \end{center}
\end{minipage}%
\hfill
\begin{minipage}[t]{0.50\textwidth}
    \vspace{0pt}
    \noindent\textbf{100.} Trên mọi đường cong hàm mũ $y = b^x$ ($b > 0, b \ne 1$), luôn tồn tại duy nhất một điểm $(x_0, y_0)$ mà tại đó tiếp tuyến của đường cong đi qua gốc tọa độ. Chứng minh rằng trong mọi trường hợp, $y_0 = e$. [\textit{Gợi ý:} Bạn có thể muốn sử dụng Công thức 1.5.10.]

    \vspace{0.2cm}
    \noindent\textbf{101.} Nếu $F = f \circ g \circ h$, trong đó $f, g$ và $h$ là các hàm khả vi, hãy sử dụng Quy tắc chuỗi để chứng minh rằng
    \[
    F'(x) = f'(g(h(x))) \cdot g'(h(x)) \cdot h'(x)
    \]

    \vspace{-0.15cm}
    \noindent\textbf{102.} Nếu $F = f \circ g$, trong đó $f$ và $g$ là các hàm khả vi hai lần, hãy sử dụng Quy tắc chuỗi và Quy tắc tích để chứng minh rằng đạo hàm cấp hai của $F$ được cho bởi
    \[
    F''(x) = f''(g(x)) \cdot [g'(x)]^2 + f'(g(x)) \cdot g''(x)
    \]
\end{minipage}

\vspace{0.35cm}

% ==============================================================================
% PHẦN GIỮA: DỰ ÁN ỨNG DỤNG (APPLIED PROJECT)
% ==============================================================================
\noindent
\begin{tcolorbox}[
    enhanced,
    colback=projectbg,
    colframe=projectbg,
    arc=0mm,
    boxrule=0pt,
    left=8pt,
    right=8pt,
    top=6pt,
    bottom=6pt
]
    % Thanh tiêu đề với các sọc ngang xanh đặc trưng chạy hai bên tiêu đề
    \noindent
    \begin{tikzpicture}[x=1cm, y=1pt]
        \foreach \y in {0, 2.3, 4.6, 6.9} {
            \draw[projectgreen, line width=0.55pt] (0, \y) -- (0.9, \y);
        }
        \node[anchor=west, inner sep=3pt] (title) at (0.9, 3.45) {
            \textsf{\textbf{\color{projectgreen}\normalsize DỰ ÁN ỨNG DỤNG}} \quad
            {\color{projectgreen}\rule[-1.5pt]{1.1pt}{10.5pt}} \quad
            \textsf{\textbf{\normalsize PHI CÔNG NÊN BẮT ĐẦU HẠ CÁNH Ở ĐÂU?}}
        };
        \foreach \y in {0, 2.3, 4.6, 6.9} {
            \draw[projectgreen, line width=0.55pt] ([xshift=3pt]title.east |- 0,\y) -- (\linewidth-18pt, \y);
        }
    \end{tikzpicture}

    \vspace{0.15cm}
    \noindent
    \begin{minipage}[t]{0.32\linewidth}
        \vspace{0pt}
        \includegraphics[width=\linewidth]{images/p0244_fig2.png}
    \end{minipage}%
    \hfill
    \begin{minipage}[t]{0.65\linewidth}
        \vspace{0pt}
        Một đường bay tiếp cận khi hạ cánh của máy bay được thể hiện trong hình vẽ và thỏa mãn các điều kiện sau:
        \begin{enumerate}[label=(\roman*),leftmargin=1.5em,itemsep=1pt,topsep=1.5pt,parsep=0pt]
            \item Độ cao bay hành trình là $h$ khi bắt đầu hạ cánh ở khoảng cách theo phương ngang $\ell$ tính từ điểm tiếp đất tại gốc tọa độ.
            \item Phi công phải duy trì vận tốc theo phương ngang không đổi $v$ trong suốt quá trình hạ cánh.
            \item Giá trị tuyệt đối của gia tốc theo phương thẳng đứng không được vượt quá một hằng số $k$ (nhỏ hơn nhiều so với gia tốc trọng trường).
        \end{enumerate}

        \vspace{0.1cm}
        \begin{enumerate}[label=\textbf{\arabic*.},leftmargin=1.5em,itemsep=2.5pt,topsep=1.5pt,parsep=0pt]
            \item Tìm một đa thức bậc ba $P(x) = ax^3 + bx^2 + cx + d$ thỏa mãn điều kiện (i) bằng cách đặt các điều kiện thích hợp cho $P(x)$ và $P'(x)$ tại thời điểm bắt đầu hạ cánh và khi tiếp đất.
            \item Sử dụng các điều kiện (ii) và (iii) để chứng minh rằng
            \[
            \frac{6hv^2}{\ell^2} \le k
            \]
            \item Giả sử một hãng hàng không quyết định không cho phép gia tốc thẳng đứng của máy bay vượt quá $k = 860\text{ mi/h}^2$. Nếu độ cao bay hành trình của máy bay là $35{,}000\text{ ft}$ và vận tốc là $300\text{ mi/h}$, phi công nên bắt đầu hạ cánh cách sân bay bao xa?
            \item[\raisebox{-0.8pt}{\begin{tikzpicture}[scale=0.85, baseline=-2.5pt]
                \draw[stewartcyan, fill=cyan!12, line width=0.5pt, rounded corners=0.5pt] (0,0) rectangle (0.42, 0.35);
                \draw[stewartcyan, line width=0.35pt] (0.21, 0) -- (0.21, 0.35);
                \draw[stewartcyan, line width=0.35pt] (0, 0.175) -- (0.42, 0.175);
                \draw[stewartcyan, line width=0.6pt] plot[smooth] coordinates {(0.05, 0.06) (0.21, 0.28) (0.37, 0.08)};
            \end{tikzpicture}}\hspace{2pt}\textbf{4.}] Vẽ đồ thị đường bay tiếp cận nếu các điều kiện nêu trong Bài toán 3 được thỏa mãn.
        \end{enumerate}
    \end{minipage}
\end{tcolorbox}

\vspace{0.45cm}

% ==============================================================================
% PHẦN DƯỚI: MỞ ĐẦU MỤC 3.5 (Implicit Differentiation)
% ==============================================================================
\noindent
\begin{minipage}[t]{0.19\textwidth}
    \vspace{0pt}
    \flushright
    \begin{tikzpicture}[baseline=(num.base)]
        \foreach \i in {0, 2.5, 5, 7.5, 10} {
            \draw[stewartcyan, line width=0.55pt] (-1.3, \i pt) -- (-0.15, \i pt);
        }
        \node[stewartcyan, font=\fontsize{26pt}{26pt}\bfseries\sffamily, anchor=base east, inner sep=0pt] (num) at (0.85, 0) {3.5};
    \end{tikzpicture}%
    \hspace{6pt}%
    {\color{stewartcyan}\rule[-4pt]{1.4pt}{24pt}}
\end{minipage}%
\hfill
\begin{minipage}[t]{0.78\textwidth}
    \vspace{0pt}
    {\bfseries\sffamily\fontsize{22pt}{24pt}\selectfont Đạo hàm ẩn}
    
    \vspace{0.25cm}
    \noindent
    {\color{stewartcyan}\rule{5pt}{5pt}}\hspace{5pt}{\bfseries\sffamily\normalsize Các hàm số xác định ẩn}
    
    \vspace{0.15cm}
    \noindent
    Các hàm số mà chúng ta đã gặp cho đến nay đều có thể được mô tả bằng cách biểu diễn một biến một cách tường minh theo biến kia---chẳng hạn,
    \[
    y = \sqrt{x^3 + 1} \qquad \text{hoặc} \qquad y = x \sin x
    \]
    hoặc, tổng quát là $y = f(x)$. Tuy nhiên, một số hàm số lại được xác định ẩn bởi một hệ thức giữa $x$ và $y$ (chẳng hạn như $x^2 + y^2 = 25$ hoặc $x^3 + y^3 = 6xy$)...
\end{minipage}

\vfill
% Dòng thông tin bản quyền chuẩn xác cuối trang sách
\noindent
\begin{minipage}{\textwidth}
    \centering
    \fontsize{6.2pt}{7.5pt}\selectfont\color{copyrightgray}
    Bản quyền thuộc sở hữu của Cengage Learning. Mọi quyền được bảo lưu. Không được sao chép, quét hoặc nhân bản toàn bộ hay một phần. WCN 02-200-203\\ 
    Do quyền điện tử, một số nội dung của bên thứ ba có thể bị ẩn khỏi sách điện tử và/hoặc chương điện tử. Đánh giá biên tập đã xác định rằng bất kỳ nội dung bị ẩn nào không ảnh hưởng đáng kể đến trải nghiệm học tập tổng thể. Cengage Learning có quyền xóa nội dung bổ sung vào bất kỳ lúc nào nếu các hạn chế quyền tiếp theo yêu cầu.
\end{minipage}

\end{document}